\documentclass[11pt]{article}
\usepackage{amsmath, amssymb}
\usepackage[margin=2.5cm]{geometry}
\title{Resep Langit \\ \large Pusaran Lautan Magma \& Air}
\author{math-aug}
\date{}

\begin{document}
\maketitle

\section{Domain Langit}

\begin{equation}
\mathcal{L}_{\text{langit}} = \{(x, y) \in \mathcal{C} \mid y_1 \le y < y_2\}
\end{equation}

dengan $y_1 = 1000$, $y_2 = 3320$, dan titik tengah $y_m = (y_1 + y_2)/2 = 2160$.

\section{Gradien Langit (Pixel Spatial Network)}

Untuk setiap piksel $(x, y) \in \mathcal{L}_{\text{langit}}$:

\begin{equation}
t(x, y) = \frac{|y - y_m|}{h_{\text{langit}} / 2}, \qquad h_{\text{langit}} = y_2 - y_1 = 2320
\end{equation}

\begin{equation}
h(x, y) = \mathrm{clip}\!\left( t(x, y) + a \sin\!\left( \frac{2\pi k x}{W} \right),\, 0,\, 1 \right)
\end{equation}

dengan $a = 0.05$, $k = 10$, $W = 7680$.

Warna RGB didefinisikan palet siang $\mathcal{P}_{\text{siang}}$
melalui interpolasi piecewise-linear:

\begin{equation}
\mathbf{C}_{\text{langit}}(x, y) = \mathcal{P}_{\text{siang}}(h(x, y))
\end{equation}

dengan anchor:
\begin{align}
\mathcal{P}_{\text{siang}}(0.0) &= (30, 80, 170) \\
\mathcal{P}_{\text{siang}}(0.5) &= (135, 190, 230) \\
\mathcal{P}_{\text{siang}}(1.0) &= (240, 248, 255)
\end{align}

\section{Awan sebagai Agregasi Elips}

Untuk $i = 0, \ldots, N - 1$ dengan $N = 500$, misal $u_i = i/(N-1)$.
Elips ke-$i$ memiliki pusat dan semi-sumbu:

\begin{align}
\mathbf{c}_i &= \bigl( W u_i,\; y_m + A_1 \sin(3\pi u_i) + A_2 \sin(7\pi u_i) \bigr) \\
s_i^x &= \tfrac{1}{2}\bigl( s_{\min} + s_{\max} \sin^{p}(\pi u_i) \bigr) \\
s_i^y &= \rho\, s_i^{x} \\
\alpha_i &= \alpha_{\min} + \alpha_{\max} \sin^{q}(\pi u_i)
\end{align}

Parameter: $A_1 = 300$, $A_2 = 200$, $s_{\min}=150$, $s_{\max}=500$,
$p=4$, $\rho=0.35$, $\alpha_{\min}=20$, $\alpha_{\max}=100$, $q=2$.

Warna elips: putih $(255, 255, 255)$ dicampur dengan warna langit
melalui compositing alpha $\alpha_i$.

\section{Peran Spatial Compression}

Eksponen $p = 4$ dan $q = 2$ menerapkan \emph{spatial compression}:
puncak $\sin^{p}(\pi u)$ memusatkan elips besar dan padat di sekitar
$u = 1/2$ (pusat horizontal kanvas), sedangkan di tepi $u \to 0$ dan
$u \to 1$ elips mengecil dan memudar secara halus. Gradien kerapatan
ini dipersepsikan mata sebagai struktur awan yang organik.

\end{document}
